% ============================================================================
% KADABRA DHT: DISTRIBUTED KNOWLEDGE ROUTING
% ============================================================================
\subsection{Kadabra: Affinity-Routed Distributed Hash Table}
\label{sec:kadabra_dht}

The Kadabra layer organizes MarkovTrie instances into a Kademlia-style
distributed hash table where nodes \emph{are} tries.  Values route to
tries based on affinity similarity, creating emergent specialization:
each trie accumulates expertise in the content domain defined by its
neighborhood in affinity space.

\subsubsection{Routing Architecture}

Each node in the network is identified by a 64-bit \texttt{NodeID} and
maintains a routing table of 64 $k$-buckets (one per bit of the ID
space), each holding up to $B = 20$ peer references.  Bucket $i$
contains peers whose XOR distance from the local node has its highest
set bit at position $i$:
\begin{equation}
  \text{bucket}(a, b)
    = \bigl\lfloor \log_2(a \oplus b) \bigr\rfloor
  \label{eq:bucket_index}
\end{equation}
This produces the standard Kademlia property that nearby nodes share
low-index buckets (many common prefix bits) and distant nodes populate
high-index buckets.

\subsubsection{Affinity-Based Value Routing}

When a Value is published to the network, its 257-bit affinity vector
(\Cref{sec:entropy_routing}) determines the target nodes.  The system
computes the XOR distance between the Value's affinity and each
candidate node's affinity, then routes the Value to the $r = 3$ closest
nodes (the replication factor).  Because the affinity vector is a
locality-sensitive hash of the token content
(\Cref{sec:capacity_analysis}), Values with similar content naturally
cluster on the same nodes.  This produces \emph{emergent specialization}:
each trie receives, and therefore learns from, a coherent subset of the
input stream.

\subsubsection{Adaptive Peer Selection via Multi-Armed Bandit}
\label{sec:mab_peers}

Peer quality within each bucket is tracked via two online statistics:
latency (exponential moving average of round-trip time) and success rate
(fraction of recent queries that returned results).  Every
$E_{\text{epoch}} = 100$ queries, the system re-evaluates each bucket:
\begin{enumerate}[nosep]
  \item Compute a quality score for each peer:
        $q = \alpha_{\text{success}} \cdot \text{success\_rate}
           - \alpha_{\text{latency}} \cdot \text{latency\_ema}$.
  \item With exploration probability $\epsilon$, probe a random
        candidate not currently in the bucket.
  \item If the candidate's quality exceeds the worst current peer,
        swap them.
\end{enumerate}
This multi-armed bandit strategy ensures that the routing table
continuously improves without requiring global coordination.

\subsubsection{Gossip Protocol and Field Digests}
\label{sec:gossip}

At each epoch boundary, every node broadcasts a compact
\texttt{FieldDigest} to its routing peers.  A digest contains:
\begin{itemize}[nosep]
  \item The originating \texttt{NodeID}.
  \item The node's 257-bit affinity vector.
  \item Surprisal mean $\bar{S}$, variance $\sigma_S^2$, and
        previous-epoch mean $\bar{S}_{\text{prev}}$ (for computing
        phase velocity).
  \item Classification entropy $H$.
  \item Node growth rate $\dot{N}$ (trie node creation rate).
  \item Effective trie depth $d_{\text{eff}}$.
  \item Epoch counter.
\end{itemize}
Digests are lightweight (a few hundred bytes) relative to the trie state
they summarize.  Each receiving node absorbs the digest into its
\texttt{FieldView}, which triggers eigenmode recomputation and field
projection (\Cref{sec:field}).  Digests are deduplicated by epoch
number: a digest with an epoch $\leq$ the stored epoch for that origin
is silently dropped.

\subsubsection{Replication and Redundancy}

Each Value is stored on $r = 3$ nodes selected by affinity proximity.
This serves two purposes: fault tolerance (any single node failure
leaves two replicas) and learning diversity (multiple tries learn from
the same data, potentially extracting different structural patterns
based on their existing knowledge).  The replication factor is
configurable but defaults to 3, consistent with standard Kademlia
deployments.
